\section{Conclusion}
\label{sec:conclusion}

In this report, we presented Douyin Multimodal Embedding (DME), a representation model for large-scale industrial multimodal retrieval. Through a two-stage training framework, DME first performs large-scale contrastive pre-training to establish a unified multimodal embedding space, and then supplements semantic sufficiency in Stage 2 via Evidence-Grounded Typed Latent Reasoning and Cross-Conditional Reconstruction, which respectively ground the embedding in retrieval-relevant evidence and enforce the preservation of fine-grained counterpart-side semantics. Cross-Conditional Reconstruction is used only during training, whereas Stage 2-A
retains a small number of latent tokens within the same encoder forward pass.
DME therefore remains a dense bi-encoder retriever, with only marginal
query-side latency overhead. In our experiments, DME achieves state-of-the-art results on MMEB-v2 among models of comparable scale, and our latency analysis confirms that the latent reasoning tokens introduce only marginal query-encoding overhead. We further quantify semantic sufficiency through a representation completeness measure, showing that counterpart-side content can be recovered from the learned embeddings at high token-level accuracy, which together with the latency analysis demonstrates that DME attains strong advantages in both semantic sufficiency and efficiency. Moreover, DME has been deployed in Douyin multimodal search, where it delivers significant gains across downstream businesses.

Looking forward, we are extending DME along two scaling directions to further strengthen the generality of the learned representations. The first is data scaling, where we enlarge and diversify the multimodal training corpus to broaden modality and task coverage. The second is model-size scaling, where we train DME with larger backbones to increase representational capacity. Our preliminary attempts along both directions already yield consistent gains, and we regard scaling data and model size as a promising path toward more general-purpose multimodal embeddings.



\newpage

\section{Contributor List}


\contributionlist

\textbf{Contributors:} \\
Haonan Chen$^{2,*,\dagger}$ \quad Chu Li$^{1,*}$ \quad Zhicheng Wang$^{1,*}$ \quad Yuanwei Liu$^{1}$ \quad Yuanjiang Wang$^{1}$

\textbf{Project Leader:} \\
Shaohua Jiang$^{1}$

\textbf{Supervisor:} \\
Zhicheng Dou$^{2}$

\textbf{Affiliations:} \\
$^{1}$ ByteDance Douyin Search Multimodal Team \\
$^{2}$ Gaoling School of Artificial Intelligence, Renmin University of China

\vspace{2pt}


\noindent
$*$ Equal contribution. Contributors are listed in alphabetical order by last name initial.


\noindent
$^\dagger$ Work was done during Haonan's internship at ByteDance Douyin Search Multimodal Team.



\section{Acknowledgments}
We sincerely thank Leyang Wang, Lanqing Hu, Tianlong Ma, Jianfeng Li and Xiangyuan Ren for their valuable support and contributions.